\begin{textsample}{POS Dim 2 – 2022\_c5 – Score -1.00 – 2022\_c5/enade\_2022\_teo\_00a9aa1edd.txt}  \label{ex:f2_pos_218}
CAPÍTULO 6 DISTRIBUIÇÃO DE CONCEITOS A finalidade deste capítulo é apresentar um panorama nacional da distribuição do Conceito Enade 2022.
Para traçar esse panorama, recorre -se a tabelas e gráficos, que são acompanhados de comentários e análises que contribuem para a contextualização dos dados neles expostos.
Para visualizar o panorama do Enade/2022 em nível nacional, na primeira seção do capítulo, os dados sobre o número de cursos participantes são organizados por grande região.
Na segunda seção, o mapeamento do conceito do exame é feito por categoria administrativa das IES e modalidade de oferta dos cursos, de maneira cruzada às grandes regiões.
Na última seção, a lógica de apresentação dos dados é feita por organização acadêmica e, também, grande região de funcionamento dos cursos, no país como um todo. 6.1 PANORAMA NACIONAL DA DISTRIBUIÇÃO DOS CONCEITOS A lógica adotada na apresentação dos dados, nesta seção e nas duas seguintes, pressupõe o registro de duas informações de natureza reiterativa.
A primeira informação diz respeito aos valores possíveis do Enade, que variam em níveis de 1 a 5, sendo que, à medida que esses valores aumentam, o curso recebe uma avaliação melhor ; além disso, os cursos com somente um ou sem nenhum concluinte participante não obtêm o Conceito Enade, ficando Sem Conceito ( SC ), como já ressaltado neste relatório.
O segundo registro é que a diferença entre os cursos tabulados neste capítulo e no Capítulo 3 corresponde aos cursos que não foram avaliados, ou seja, àqueles, em princípio, sem estudantes concluintes que tivessem sido inscritos para a prova.
Na Tabela 6.1, são apresentadas a quantidade e a distribuição de cursos de Teologia participantes do Enade/2022, por faixa de conceito e grande região.
Observando -se os dados da Tabela 6.1, nota -se que, dos 131 cursos participantes, 51 ( 38,9 \% ) classificaram -se com conceito 3, o valor modal.
Esse foi também o conceito modal na maioria das regiões : Norte ( 33,3 \% ), Sudeste ( 42,6 \% ) e Sul ( 41,7 \% ).
O conceito 2 foi o segundo mais frequente em nível nacional ( 25,2 \%, correspondendo a 33 cursos ), tendo sido o conceito modal na região Centro-Oeste ( 46,2 \% ) e Nordeste ( 36,4 \% ).
O conceito 4 foi o terceiro ( 19,1 \%, correspondendo a 25 cursos ).
Houve, ainda, 3 cursos ( 2,3 \% ) que receberam o conceito 5, e 9 cursos ( 6,9 \% ) que receberam o conceito 1.
Dos 131 cursos de Teologia, 10 ( 7,6 \% ) ficaram Sem Conceito ( SC ).
A região Norte participou com 6 cursos, o que corresponde a 4,5 \% do total nacional.
Desses, 2 cursos receberam o conceito 2 e 3, cada, o que equivale a 33,3 \%, de cada um dos conceitos, do total regional.
A um curso ( 16,7 \% ) atribuiu -se o conceito 1, e a 1 curso ( 16,7 \% ), o conceito 4.
Nenhum curso ficou Sem Conceito ( SC ) ou recebeu conceito 1 e 5, nessa região.
A região Nordeste participou com 22 cursos, o correspondente a 16,7 \% do total nacional.
Nessa região, 8 cursos ( 36,4 \% em termos regionais ) obtiveram conceito 2, o conceito modal para a região, como já comentado.
Os conceitos 1 e 3 foram atribuídos, respectivamente, a 5 e a 7 cursos ( 22,7 \% e 31,8 \%, respectivamente ).
Ainda nessa região, 1 curso ( 4,5 \% ) foi avaliado com o conceito 4 e nenhum curso foi avaliado com o conceito 5.
Um curso ficou Sem Conceito ( SC ).
Para os 54 cursos participantes da região Sudeste ( 41,2 \% do total nacional ), o conceito 3 foi o conceito modal, como já anteriormente assinalado, tendo sido obtido por 23 cursos ( 42,6 \% ).
O conceito 1 foi atribuído a três cursos ( 5,6 \% ), e o conceito 2, a 11 cursos ( 20,4 \% ).
Receberam os conceitos 4 e 5, respectivamente, 12 cursos ( 22,2 \% ) e 1 curso ( 1,9 \% ).
Nessa região, quatro cursos ficaram Sem Conceito ( SC ).
Os 36 cursos da região Sul corresponderam a 27,4 \% do total nacional.
Como já indicado, houve predominância do conceito 3 ( conceito modal ), atribuído a 41,7 \% dos cursos da região, ou seja, atribuído a 15 dos 36 cursos participantes na região Sul.
O conceito 1 foi atribuído a nenhum curso ; o conceito 2, a 6 cursos ( 16,7 \% ) ; e os conceitos 4 e 5, a 9 ( 25 \% ) e a dois cursos ( 5,6 \% ), respectivamente.
Na região Sul, três cursos ficaram Sem Conceito ( SC ).
Já, dos 13 cursos participantes na região Centro-Oeste ( 9,9 \% do total nacional ), 6 ( 46,2 \% em termos regionais ) receberam conceito 2, o conceito modal.
Nenhum curso obteve o conceito 1, 4 cursos ( 30,8 \% ) obtiveram o conceito 3, 2 ( 15,4 \% ), o conceito 4, e nenhum, o conceito 5.
Nessa região, tivemos um curso que ficou Sem Conceito ( SC ).
Em um gráfico de distribuição cumulativa, em cada valor de abscissa ( eixo x ), o valor apresentado no eixo das ordenadas ( eixo y ) é o valor acumulado da variável ( Conceito Enade ).
Em linhas gerais, em um gráfico dessa mesma tipologia, poligonais mais à esquerda representam, no caso dos dados apresentados nesta seção, as grandes regiões com a pior distribuição, e poligonais mais à direita, as grandes regiões com os melhores conceitos.
A partir dessas diretrizes, no Gráfico 6.1, é apresentada a distribuição cumulativa do Conceito Enade, segundo a grande região.
Nesse gráfico, poligonais mais à esquerda representam valores mais baixos do Conceito Enade e poligonais mais à direita, valores mais altos. 6.2 CONCEITOS POR CATEGORIA ADMINISTRATIVA, MODALIDADE DE OFERTA DOS CURSOS E GRANDE REGIÃO Na Tabela 6.2, é apresentada a distribuição dos cursos de Teologia participantes do Enade/2022, por categoria administrativa e modalidade de oferta, de acordo com os conceitos por eles alcançados, segundo as grandes regiões brasileiras.
De acordo com as informações da Tabela 6.2, dos 131 cursos participantes, nenhum era ministrado em Instituições Públicas, e 131 ( 100 \% ), em Instituições Privadas.
Quanto à modalidade de oferta, 49 cursos eram a Distância, e 82, presenciais.
Ainda, conforme as informações da mesma tabela, em termos nacionais, observa -se que apenas 3 cursos foram avaliados com conceito 5.
Na rede privada, o conceito modal foi o 3, atribuído a 51 cursos, o correspondente a 38,9 \% dos 131 cursos da categoria.
Dos demais cursos participantes, 7 ( 5,3 \% ) receberam o conceito 1, e 37 ( 28,2 \% ), o conceito 2.
O conceito 4 foi atribuído a 24 cursos ( 18,3 \% ).
Nessa categoria administrativa, 9 cursos ( 6,8 \% ) ficaram Sem Conceito ( SC ).
Dos 82 cursos participantes da modalidade presencial, o conceito 3 foi o valor modal, atribuído a 35 cursos.
Dos demais cursos nessa modalidade de oferta, 7 cursos receberam o conceito 1 ; 13 cursos receberam o conceito 2 ; 22 cursos, o conceito 4 ; e 3 cursos, o conceito 5.
Nessa modalidade de oferta, dois cursos ficaram Sem Conceito ( SC ).
Dos 49 cursos participantes da modalidade a distância, o conceito 2 foi o valor modal, atribuído a 20 cursos.
Dos demais cursos nessa modalidade de oferta, 2 cursos receberam o conceito 1 ; 16 cursos receberam o conceito 3 ; 3 cursos, o conceito 4 ; e nenhum curso, o conceito 5.
Nessa modalidade de oferta, oito cursos ficaram Sem Conceito ( SC ).
Na análise por região, observa -se que, na região Norte, as instituições privadas participaram com 6 cursos ( 4,5 \% do total regional ), e nenhum curso ficou Sem Conceito ( SC ).
O conceito modal para as instituições privadas na região foi o 2 e 3, atribuído a 2 cursos, cada.
Um curso foi avaliado com o conceito 1 ; 1, com o conceito 4 ; e a nenhum curso foi atribuído o conceito 5.
Dos cursos oferecidos na região Norte, 4 eram presenciais, e nenhum ficou Sem conceito ( SC ).
Dos 4 cursos oferecidos, dois obtiveram o conceito 2 ; 1, o conceito 3 ; um, o conceito 4 ; e nenhum o conceito 5.
Dois cursos eram oferecidos na modalidade a distância, sendo que um recebeu o conceito 1 ; e 1, o conceito 3.
Nenhum curso ficou Sem Conceito ( SC ) ou recebeu os demais conceitos.
Na região Nordeste, dentre os 22 cursos oferecidos por IES privadas, 8 cursos foram avaliados com conceito 2, o conceito modal.
Cinco cursos obtiveram o conceito 1 ; sete, o conceito 3 ; 1, o conceito 4 ; e nenhum curso obteve o conceito 5.
Nessa categoria administrativa e grande região, um curso ficou Sem Conceito ( SC ).
Dos cursos oferecidos na região Nordeste, 16 eram presenciais ; desses, nenhum ficou Sem Conceito ( SC ), cinco ficaram com o conceito 1 ; 3, com o conceito 2 ; 7, com o conceito 3, o valor modal ; 1, com o conceito 4 ; e nenhum, com o conceito 5.
Na modalidade a distância, eram 6 cursos, tendo nehum obtido o conceito 1 ; 5, o conceito 2 ; nenhum, o conceito 3, 4 ou 5.
Nessa região, um curso ficou Sem Conceito ( SC ).
Na região Sudeste, temos um total de 54 cursos em IES privadas, sendo o conceito modal o 3, atribuído a 23 cursos.
Quatro cursos ficaram Sem Conceito ( SC ).
Os demais foram avaliados com os seguintes conceitos : 1 ( três cursos ), 2 ( 11 cursos ), 4 ( 12 cursos ) e 5 ( 1 curso ).
Também na região Sudeste, grande parcela dos cursos foi oferecida na modalidade presencial : 32 dos 54 cursos.
Para esses, a moda foi o conceito 3, com 15 cursos.
Os demais cursos presenciais da região foram alocados nos conceitos 2 ( 3 cursos ), 4 ( 11 cursos ) e 5 ( 1 curso ).
Nenhum curso ficou Sem Conceito ( SC ) na modalidade presencial na região Sudeste.
Na região Sudeste, 22 cursos foram oferecidos na modalidade a distância : um curso com conceito 1 ; 8 com conceito 2 ; 8 com conceito 3 ; 1 com o conceito 4 ; e nenhum com o conceito 5.
Na região Sudeste, quatro cursos a distância ficaram Sem Conceito ( SC ).
As instituições privadas concentraram 36 cursos participantes da região Sul.
Desses, 15 cursos obtiveram conceito 3, o conceito modal.
Nessa combinação de categoria administrativa e grande região, nenhum curso recebeu o conceito 1 ; seis receberam o conceito 2 ; outros nove, o conceito 4 ; e ainda dois, o conceito 5.
Quatro cursos ficaram Sem Conceito ( SC ).
Na região Sul, 21 dos cursos foram oferecidos na modalidade presencial.
Para esses, a moda foi o conceito 3, com 9 cursos.
Os demais cursos presenciais da região foram alocados nos conceitos 2 ( 2 cursos ), 4 ( 7 cursos ) e 5 ( 2 cursos ).
Nenhum curso ficou Sem Conceito ( SC ) nesta região.
Na modalidade a distância, foram oferecidos 15 cursos, sendo que nenhum curso recebeu conceito 1 ; 4 ficaram com conceito 2 ; 6 com conceito 3 ; e 2 com o conceito 4.
Três cursos ficaram Sem Conceito ( SC ) e nenhum curso ficou com o conceito 5.
Na região Centro Oeste, tivemos 13 cursos participantes de IES privadas.
Desses, 6 cursos receberam conceito 2, o valor modal.
Dos demais cursos, nenhum recebeu o conceito 1 ; 4 receberam o conceito 2 ; 2, o conceito 4 ; e nenhum curso recebeu o conceito 5.
Um curso ficou Sem Conceito ( SC ).
Também na região Centro Oeste, a maioria dos cursos foi presencial : 9 ( 69,2 \% ).
Para esses, a moda foi o conceito 2 e 3, com 3 cursos cada ( 33,3 \% ).
Os demais cursos presenciais da região foram alocados nos conceitos 4 ( 2 cursos, 9,0 \% ).
Um curso presencial da região ficou Sem Conceito ( SC ) e nenhum recebeu o conceito 1 e 5, nessa modalidade.
Para os 4 cursos a distância ( 30,7 \% ), a moda foi para o conceito 2 ( 3 cursos ).
Os demais cursos a distância da região receberam os conceitos 3 ( um curso ).
Nenhum curso a distância da região ficou Sem Conceito ( SC ) ou com o conceito 1, 4 e 5.
À título de síntese, no Gráfico 6.2, apresenta -se a distribuição do Conceito Enade, segundo a categoria administrativa da IES. 6.3 CONCEITOS POR ORGANIZAÇÃO ACADÊMICA E GRANDE REGIÃO Na Tabela 6.3, encontra -se a distribuição dos conceitos atribuídos aos cursos participantes do Enade/2022, na área de Teologia, por organização acadêmica, segundo as grandes regiões brasileiras.
Como mostram os dados da Tabela 6.3, dos 131 cursos de Teologia participantes, 32 eram oferecidos em universidades, 23 em centros universitários, 76 em faculdades e nenhum em CEFET/IF.
Esta distribuição corresponde a, respectivamente, 24,4 \%, 17,5 \% e 58 \% dos cursos.
De acordo com os dados apresentados na mesma tabela, de todos os 3 cursos avaliados com o conceito 5, 1 era vinculado a universidades.
Os cursos vinculados a esse tipo de organização acadêmica tiveram o conceito 3 como conceito modal, atribuído a 15 cursos ( 46,8 \% ).
Os demais cursos ofertados por universidades receberam os conceitos 2 ( 8 cursos ), 4 ( 5 cursos ) e conceito 5 ( 1 curso, como já mencionado ).
Três cursos ficaram Sem Conceito ( SC ).
Dos cursos vinculados a centros universitários, o conceito modal também foi o 3, atribuído a 8 cursos ( 34,7 \% ).
Cinco cursos vinculados a esse tipo de organização acadêmica ficaram Sem Conceito ( SC ), e os outros cursos receberam os conceitos 1 ( 1 curso ), 2 ( 7 cursos ) e 4 ( 2 cursos ).
Dos 76 cursos mantidos por faculdades, 28 ( 36,8 \% ) receberam conceito 3, o conceito modal.
O segundo conceito mais frequentemente atribuído aos cursos oferecidos por essa categoria de IES foi o conceito 2, recebido por 22 cursos.
Na sequência, 17 cursos obtiveram conceito 4 ; 6 cursos, o conceito 1 ; e dois cursos, o conceito 5.
Dos cursos oferecidos em faculdades, um ficou Sem Conceito ( SC ).
Considerando -se, separadamente, as regiões brasileiras, verifica -se que, na região Norte, as universidades participaram com 6 cursos avaliados, aos quais foram atribuídos os seguintes conceitos : 2, a 2 cursos ; 3, a 2 cursos ; 4 a um curso ; 1, a um curso ; e 5, a nenhum curso.
Nenhum curso ficou Sem Conceito ( SC ) nessa combinação de categoria e região.
As universidades da região Norte foram representadas por 1 curso, que obteve conceito 3, o valor modal.
Os conceitos 1, 2, 4 e 5 não foram atribuídos a nenhum curso.
Do mesmo modo, nenhum curso ficou Sem Conceito ( SC ).
As faculdades participaram com 5 cursos na região Norte, os quais receberam os conceitos 1 ( 1 curso ), 2 ( dois cursos, conceito modal ), 3 ( 1 curso ) e 4 ( um curso ).
Nenhum curso ficou Sem Conceito ( SC ), e nenhum curso oferecido em faculdades da região obteve o conceito 5.
Nessa região e organização acadêmica, nenhum curso foi oferecido em CEFET/IF ou centro universitário.
Na região Nordeste, as universidades participaram com 3 dos 22 cursos da área de Teologia oferecidos.
O conceito modal foi 3, atribuído a 2 cursos.
Um curso obteve o conceito 2 ; nenhum curso recebeu o conceito 1, 4 e 5.
Nenhum curso oferecido por universidades no Nordeste ficou Sem Conceito ( SC ).
Os centros universitários contaram com 3 cursos participantes na região Nordeste, dos quais os conceitos 2 e 4 foram dados a um curso, cada.
Os outros conceitos não foram dados a nenhum outro curso.
Um curso oferecido por centros universitários no Nordeste ficou Sem Conceito ( SC ).
As faculdades foram representadas por 16 cursos na região Nordeste, 7 dos quais receberam conceito 2, o conceito modal.
Aos demais cursos foram atribuídos os conceitos 1 ( 4 cursos ), 3 ( 5 cursos ).
O conceito 4 e 5 não foram dados a nenhum curso.
Nenhum curso ficou Sem Conceito ( SC ).
Nessa combinação de região e organização acadêmica, nenhum curso foi oferecido em CEFET/IF.
Na região Sudeste, as universidades concentraram 15 dos 54 cursos de Teologia da região.
Dos cursos oferecidos em IES com este tipo de organização acadêmica, nessa região, o conceito modal foi 3, atribuído a 7 cursos.
Os demais cursos receberam os conceitos 2 ( 3 cursos ) e 4 ( 4 cursos ).
Dois cursos ficaram Sem Conceito ( SC ), e nenhum recebeu o conceito 1 e 5.
Os centros universitários participaram com 11 cursos na região Sudeste, dos quais 5 obtiveram o conceito modal 3.
O conceito 1 foi atribuído a um curso ; o conceito 2, a 3 cursos ; o conceito 4, a 1 curso ; e o conceito 5, a nenhum curso.
Um curso ficou Sem Conceito ( SC ) nessa combinação de organização acadêmica e região.
Já as faculdades foram representadas por 27 cursos na região Sudeste, dos quais um ficou Sem Conceito ( SC ).
O conceito modal 3 foi atribuído a 11 cursos.
Aos demais cursos foram atribuídos os conceitos 1 ( 0 cursos ), 2 ( 7 cursos ), 4 ( 7 cursos ) e 5 ( um curso ).
Dos 36 cursos da região Sul, 8 eram mantidos por universidades, para os quais o conceito modal foi 3, atribuído a 4 cursos.
Os demais cursos receberam os conceitos 2 ( 2 cursos ) e 5 ( 1 curso ) ; nenhum curso recebeu o conceito 1 ou 4.
Em IES com esse tipo de organização acadêmica, um curso da região Sul ficou Sem Conceito ( SC ).
Dos 7 cursos participantes de centros universitários da região Sul, dois ficaram Sem Conceito ( SC ), e 3 cursos receberam conceito 3, o valor modal.
Os conceitos restantes ficaram assim distribuídos : 2 cursos ( conceito 2 ) ; nenhum curso recebeu o conceito 1, 4 ou 5.
Quanto aos 21 cursos vinculados a faculdades na região Sul, 8 cursos receberam o conceito 3, valor modal em conjunto com o conceito 4, que foi dado a 8 cursos.
Nenhum curso ficou Sem Conceito ( SC ).
Os demais cursos receberam os conceitos 1 ( um curso ), 2 ( 3 cursos ), e 5 ( um curso ).
Na região Centro-Oeste, 57 dos 137 cursos oferecidos eram mantidos por universidades.
Para esse tipo de organização acadêmica, o conceito modal foi 3, atribuído a 28 cursos ; e, nenhum curso da região ficou Sem Conceito ( SC ).
Os demais cursos receberam os conceitos 1 ( três cursos ), 2 ( 16 cursos ), 4 ( 7 cursos ) e 5 ( três cursos ).
Os centros universitários da região Centro-Oeste contaram com 2 cursos : 1 curso foi avaliado com o conceito 2.
Não houve cursos avaliados com o conceito 1, 3, 4 e 5 e um curso ficou Sem Conceito ( SC ).
Dos 7 cursos oferecidos em faculdades na região Centro-Oeste, três cursos receberam o conceito 2 ; 3 cursos, o conceito 3 ; 1 curso, o conceito 4.
Nessa região, nenhum curso oferecido em faculdades ficou Sem Conceito ( SC ) ou obteve o conceito 1 ou 5.
De maneira semelhante à forma como se procedeu, em relação à seção anterior, encerra -se esta com o Gráfico 6.3, que apresenta a distribuição do Conceito Enade dos cursos de Teologia segundo a organização acadêmica da IES.
No Gráfico 6.3, apresenta -se a distribuição do Conceito Enade segundo a organização acadêmica da IES.
A leitura desse último gráfico é semelhante aos anteriores, seguindo -se a lógica de que poligonais deslocadas para a esquerda correspondem avalores mais baixos do Conceito Enade ( 1, 2 e 3 ) e poligonais deslocadas para a direita correspondem a valores mais altos do Conceito Enade.

% matched lemmas: 
\end{textsample}
